Suppose you are travelling on a spacecraft that encounters a uniformly spherically distributed gas cloud with radius $R$, and volume density $\rho$. At your moment of approach, you decelerate till you are at rest with respect to the gas cloud and then turn off your engines. You may neglect gas friction in this problem.
\textbf{Note.} Gravitational potential inside the uniform solid sphere is given by
\[V(r)=-\frac{GM}{R^3}\left(\frac{3R^2}{2}-\frac{r^2}{2}\right),\]
where $r$ is the distance from the center of the gas cloud. You may also find the following fact useful. For an equation of the form $a=-k^2r$, where $k$ is some constant and $a$ is acceleration, the time period of the oscillation is given by $T=2\pi/k$.
\begin{enumerate}
\item Once inside the gas cloud, how long will it take you to return to your original position if you are in free fall inside the cloud?
\item Start from the law of conservation of energy. Suppose you are tired of being inside the gas cloud after investigating it for a while and want to escape it (to infinity). What is your required speed if you are initially located in the center of the cloud and are at rest?
\end{enumerate}
